\section{2026 年 11 月核验增补}

本节按论文首次提交时间排列。每项方法附有可核验的外部入选依据。

% METHOD: 2312.02189
\paragraph{StableDreamer：分阶段控制文本到高斯生成的 SDS 噪声}
2023.12~\cite{2312_02189}。StableDreamer 将分数蒸馏采样的更新改写成对去噪图像伪目标的监督，以此观察不同噪声水平怎样影响三维优化。它用图像空间扩散先改善粗几何与多面伪影，再切换到潜空间扩散细化高斯颜色和纹理，两个阶段采用不同的退火策略。

论文以同类提示词对比 DreamFusion、Magic3D、GSGen 和 ProlificDreamer 的几何、外观及训练和渲染开销。对照主要展示生成样例和效率；渲染可超过每秒 30 帧描述的是已优化高斯，不代表单次文本生成耗时，也不意味着各种提示的多面伪影都能消除。论文作者页列出 Apple。

\textbf{入选依据} 论文署名单位包括 Apple（\href{https://arxiv.org/abs/2312.02189}{论文原文}）。

% METHOD: 2312.09682
\paragraph{DreamGaussian 的仰角敏感性：单图动物资产小样本观察}
2023.12~\cite{2312_09682}。这项研究复用 DreamGaussian 生成动物网格，将估计或指定的输入相机仰角作为变化条件，观察参考方向如何影响高斯优化及后续网格。研究并未引入新的高斯表示或训练损失。

原文用棕熊、鹳和蜻蜓各一幅图，比较负 30 至正 30 度的七档仰角。图示说明不合适的仰角会改变形体与多视图外观，但样本数量小，主要是定性观察，没有三维参考测量来检验生物结构精度。作者单位列为华为土耳其研发中心。

\textbf{入选依据} 论文署名单位包括 Huawei（\href{https://arxiv.org/abs/2312.09682}{论文原文}）。

% METHOD: 2608.17988
\paragraph{GS-Voxel：为大范围航拍场景构造结构化高斯潜变量}
2026.08~\cite{2608_17988}。GS-Voxel 将预优化且满足输入约定的高斯集合确定性组织为稀疏体素。因子化 VAE 分别编码体素占据和局部高斯属性，图像条件流模型在潜空间生成，带重叠的平铺推理将合成范围拓展到单块训练裁剪之外。这里避免的是逐场景向固定大小模板拟合，训练资料仍来自已重建的 3DGS 场景。

论文在各含 1.5 万张渲染图像的生成与留出参照集合上报告 FID 28.0、KID 0.020，并明确区分不同论文的相机协议和数据。当前实现只在 SH0 航拍场景开发和验证，改变卫星影像来源后的鲁棒性尚未评估。极密体素的容量限制和有限空间分辨率也可能分别造成原语丢弃与细结构退化。

\textbf{入选依据} 论文署名单位包括 Alibaba（\href{https://arxiv.org/abs/2608.17988}{论文原文}）。

% METHOD: 2609.03931
\paragraph{SPAR3S：在稀疏体素潜空间自回归补全场景}
2026.09~\cite{2609_03931}。SPAR3S 将稀疏条件图像编码到与高斯解码器对齐的体素潜空间。掩码自回归模型在其中联合预测新增占据和潜变量，随后解码出可供新视点渲染的高斯集合。研究针对条件视图很少、彼此甚至没有重叠的室内场景重建。

论文在合成室内 3DFront 与真实视频 RealEstate10K 上评估，前者提供真值位姿并从渲染深度回投点图，后者由 MASt3R-SfM 估计初始几何。作者指出占据预测漏掉墙壁或物体体素时会留下空洞，有限体素分辨率和每体素的高斯密度也会引入局部模糊。

\textbf{入选依据} 论文署名单位包括 NAVER LABS（\href{https://arxiv.org/abs/2609.03931}{论文原文}）。

% METHOD: 2610.23050
\paragraph{TriFold：用规则三平面编码不规则高斯场景}
2026.10~\cite{2610_23050}。TriFold 将带有空间密度差异和可变原语数量的场景高斯集合编码为固定布局的三平面变分潜变量。编码采用单元查询及以原语为锚的跨平面注意力，解码器再按占据与局部原语置信预测分配高斯。规则潜空间也供条件流匹配生成使用。

论文在 InteriorGS 与 SceneSplat-49K 的 Aria Synthetic Environments 子集验证重建，按原场景划分训练与测试以避免裁剪泄漏。论文另对点云条件和单图全形状补全展示生成结果，并指出观察稀疏的场景边界仍可能较真值模糊。

\textbf{入选依据} 论文署名单位包括 Alibaba（\href{https://arxiv.org/abs/2610.23050}{论文原文}）。

% METHOD: 2611.25338
\paragraph{VideoRNG：用标准化视频观测生成可重照明的神经高斯}
2026.11~\cite{2611_25338}。VideoRNG 在规定的转台相机与照明下从文本或图像合成观察序列，将 29 帧编码至 8 个视频潜帧，前景特征经重建网络一次预测高斯几何与神经外观。几何与外观的组织允许在新观察方向和新光照下渲染资产。

作者用约 2.2 万个筛选后的 TexVerse 纹理资产训练，并在留出的视点和灯光条件下评估重建。论文深度锚定对比的 PSNR 为 29.22 dB，联合深度监督为 25.39 dB，其含义限于该比较。生成视频的跨帧不一致可在重建中留下重影，尖锐镜面效果仍难复现。

\textbf{入选依据} 论文署名单位包括 NVIDIA（\href{https://arxiv.org/abs/2611.25338}{论文原文}）。
